\documentclass[12pt]{article}

\usepackage{float}
\usepackage{graphicx} % Required for inserting images
\usepackage{fontspec}
\usepackage{apacite}
\setmainfont{Times New Roman}
\usepackage[letterpaper,top=2cm,bottom=2cm,left=2cm,right=2cm,marginparwidth=1.5cm]{geometry}

\title{Weekly assignment 3: Compliance implementation}
\author{Truong Khang Tobias Nguyen}
\date{}

\begin{document}

\maketitle

\section{Deciding on a Standard}
When deciding which standard an organization should follow, the organization should ask itself which standards fulfills their specific security requirements.
\newline
Rather an just following a standard for the sake of following a standard, organizations should seek which standards suits their operations and requirements the most and whether they need it due to regulatory requirements. As standards are usually made to apply to as many organizations from different sectors as possible, organizations should also look on how they can adapt a standard for themselves specifically.
\newline
Additionally, organizations also have to evaluate whether adopting a standard will actually reduce their risks or provide them a competitive advantage. In that sense, a reasonable approach would be to first implement mandatory standards dictated by law and industry, and then including voluntary or best-practice standards that align with an organization's own risk tolerance and strategic goals.
\newline
To conclude, organization's should choose standards based on whether they are mandated or not and how they align with the organization itself.

\section{Gaps in adopting Standards}
Despite the vast amount of standards available that organizations can adopt, the reality is that there a still organizations that don't adopt standards such as Vaastamo, where failure occurred despite established practices existing.
\newline
I believe that this gap is not caused by a lack of knowledge but rather from organizational, economic and cultural barriers. Firstly, Costs and complexity can be a major deterrent as implementing a standard usually require investments in technology, training, personnel and consulting. These costs may lead to organizations rather pursuing short-term productivity and financial goals.
\newline
Secondly, organizational issues can play another role why organizations don't adopt standards. As security isn't generating revenue and isn't immediately producing visible results, it can be seen as a cost rather than an investment which subsequently leads to organizations not pursuing standards. Additionally, organizations can look at security measures as general IT problems rather than disruptions that can cause downtimes in their operations.
\newline
Finally, there are cultural challenges as if compliance is seen as optional or poorly enforced in organizations, they may not feel the need to implement standards. If organizations prioritize operations over establishing a strong cybersecurity culture, they allow non-compliant behaviour to be an acceptable risk. 
\newline
Picking up the Vaastamo case again, one assumption on why they did not have any standards is a combination of all the aforementioned factors. 
\newline
From a financial standpoint, even though Vaastamo was doing financially well, their IT department, which given that the CEO blamed them for their data breach and thus can be assumed to have had the responsibility for the organization's security, had virtually no budget. Therefore, even if they wanted to implement a standard, they did not have the resources for the necessary tools, training and support. 
\newline 
Organizationally, as the CEO claimed that the IT department was at fault, the assumption is that security was seen as an IT responsibility rather than a strategic goal and thus, they might have seen implementing a standard as not necessary.
\newline
From a cultural standpoint, Vaastamo had several vulnerabilities that were not addressed prior to their incident. In that sense, I argue that Vaastamo saw security measures as optional and subsequently also viewed standards as optional.

\section{Potential Improvements}
As organizations ultimately decide for themselves whether to follow standards or not, the only way I can think of is to provide incentives on why organizations should adopt standards.
\newline
Organizations should know why standards exit and how they can be used to minimize risks. Additionally, organizations should also be made aware of the consequences that can occur when they do not follow any standards, e.g. Vaastamo.
\newline
But in order to make organizations follow standards without enforcing them, requires changes in the standards themselves. Standards lack flexibility as they seek to apply to as many organizations as possible. Due to that, some measures of standard may apply to one organization but not to another. In that sense, they should allow a certain degree of freedom without it negatively influencing audits or the chances of a certification.
\newline
Additionally, as standards have the risk of becoming outdated, standards need to continuously monitor trends and technology to stay up to date and this should also not come to a disadvantage to organizations seeking to adopt standards.
\newline
The brute force approach to make organizations adopt standards is to enforce them through legislation. Even though every organization within a system will follow the general baseline, I believe that it will reinforce my two points about changing the nature of standards. Forcing organizations to implement measures that they do not need because it is mandated and the risk of becoming outdated may lead to loss of acceptance of standards themselves.
\newline
Weak acceptance in combination with legislative enforcement can lead to organizations adopting standards as a bare minimum measure without adaptions to their specific operations.
This is problematic as standards should be seen as a layer in an organization's security strategy that should be complemented with additional specific measures rather than the entire strategy.
\newline
To conclude, I believe that the adoption of standards should not be enforced through legislative measures but organizations should be incentivized or persuaded to adopt standards.

\end{document} 